\documentclass[12pt]{article}
\usepackage[utf8]{inputenc}
\usepackage[T1]{fontenc}
\usepackage{hyperref}
\title{Far Western Muria Language Resources for K-12 Curriculum Development}
\author{Compiled by Claude Sonnet 5.5 (claude-sonnet-5-5) for Andr\'as Kornai}
\date{}
\begin{document}
\maketitle

\section{Introduction}
Far Western Muria (also Gaita Koitor) is a Dravidian language of the Gondi group, spoken by Muria Adivasi communities in the Bastar region of Chhattisgarh, India. It is one of three Muria varieties (with Eastern and Western Muria). It is sometimes confused with Madiya. Very little was found in the searches.

\section{Primary Educational Resources}
None found.

\section{Linguistic and Reference Materials}
\begin{itemize}
\item Muria language (Wikipedia): \url{https://en.wikipedia.org/wiki/Muria_language}
\item Garapati, Umamaheshwar Rao 1991, study of Gondi dialect subgrouping that includes FW Muria (Glottolog reference record): \url{https://glottolog.org/resource/reference/id/475796}
\end{itemize}

\section{Religious and Community Resources}
None found.

\section{Digital and Audio Resources}
None found.

\section{Curriculum Development Context}
The language is not known to be taught in schools, and no orthography or teaching materials were found. Any curriculum would require basic documentation and community work first.
\end{document}